	\begin{UseCase}{CU4}{Gestionar Proyectos de Colaborador}{
		Una vez que el \hyperlink{admin}{Administrador} ha registrado un proyecto y ha asignado a un \hyperlink{jefe}{Líder}, este último (desde su cuenta), tendrá la facultad de realizar operaciones sobre el proyecto creado como la de asignar a los \hyperlink{analista}{Analistas} con los que trabajará en conjunto durante el desarrollo del proyecto. Al momento de asignar a los analistas, estos también podrán realizar ciertas operaciones sobre el proyecto.\\
		
		Este caso de uso permite al colaborador del proyecto (Líder o Analista), visualizar en una tabla el listado de los proyectos en los que se encuentra participando, así como realizar las acciones necesarias para controlar aquellos que han sido previamente registrados por el \hyperlink{admin}{Administrador} (al acceder a un proyecto le permitirá gestionar: módulos del proyecto, casos de uso, \hyperlink{tElemento}{elementos} del caso de uso y descargar el documento de análisis).\\
		
		Para el caso específico del lider se incluyen las siguientes acciones: revisar y realizar observaciones de los \hyperlink{tElemento}{elementos} para enviarlos a corrección, así como liberar los casos de uso que considere completos y correctos. Para el caso del analista: corregir la información de los \hyperlink{tElemento}{elementos} con base en las observaciones realizadas por el líder. \\
	}
	\UCitem{Actor}{\hyperlink{jefe}{Líder de análisis}, \hyperlink{analista}{Analista}}
	\UCitem{Propósito}{Visualizar los proyectos a los que el colabordor en sesión se encuentra asociado, así como entrar a cada uno de ellos para realizar las actividades correspondientes}
	\UCitem{Entradas}{Ninguna}
	\UCitem{Salidas}{\begin{itemize}
			\item  \cdtRef{proyectoEntidad}{Proyecto}: Tabla que muestra \cdtRef{proyectoEntidad:claveProyecto}{Clave}, \cdtRef{proyectoEntidad:nombreProyecto}{Nombre} y el \cdtRef{proyectoEntidad:liderProyecto}{Líder del Proyecto} de todos los registros de los proyectos.
			\item  \cdtIdRef{MSG2}{No existe información}: Se muestra en la pantalla \IUref{IU5A}{Gestionar proyectos de colaborador} cuando el actor no se encuentra asociado a ningún proyecto.
	\end{itemize}}
	\UCitem{Precondiciones}{Que el colaborador en sesión se encuentre asociado al menos a un proyecto.}
	\UCitem{Postcondiciones}{Ninguna}
	\UCitem{Errores}{Ninguno}
	\UCitem{Tipo}{Caso de uso primario}
\end{UseCase}
%--------------------------------------
\begin{UCtrayectoria}
	\UCpaso[\UCactor] Solicita gestionar los proyectos presionando la opción ''Proyectos'' del menú \IUref{MN2}{Menú de Colaborador}.
	\UCpaso[\UCsist] Obtiene la información de los proyectos en los que el actor se encuentra colaborando. \hyperlink{CU4:TAA}{[Trayectoria A]}
	\UCpaso[\UCsist] Ordena los proyectos alfabéticamente basandonse en la clave del mismos.
	\UCpaso[\UCsist] Muestra la información de los proyectos en la pantalla \IUref{IU5}{Gestionar Proyectos de colaborador}. \label{CU4-P3}
	\UCpaso[\UCsist] Muestra el botón \raisebox{-1mm}{\includegraphics[height=11pt]{images/Iconos/Colaboradores}} para cada proyecto en el que el actor sea líder. \hyperlink{CU4:TAB}{[Trayectoria B]}
	\UCpaso[\UCactor] Gestiona los proyectos a través de los botones: \raisebox{-1mm}{\includegraphics[height=11pt]{images/Iconos/Colaboradores}}, \raisebox{-1mm}{\includegraphics[height=11pt]{images/Iconos/entrar}}, \raisebox{-1mm}{\includegraphics[height=11pt]{images/Iconos/pp}} y \raisebox{-1mm}{\includegraphics[height=11pt]{images/Iconos/word}}. \label{CU4-P5}
\end{UCtrayectoria}		
%--------------------------------------
\hypertarget{CU4:TAA}{\textbf{Trayectoria alternativa A}}\\
\noindent \textbf{Condición:} El actor no se encuentra colaborando en ningún proyecto.
\begin{enumerate}
	\UCpaso[\UCsist] Muestra el mensaje \cdtIdRef{MSG2}{No existe información} en la pantalla \IUref{IU5A}{Gestionar Proyectos de Colaborador} para indicar que no hay registros de proyectos para mostrar.
	\item[- -] - - {\em {Fin del caso de uso}}.%
\end{enumerate}
%--------------------------------------


\hypertarget{CU4:TAB}{\textbf{Trayectoria alternativa B}}\\
\noindent \textbf{Condición:} El actor inicio sesión como analista.
\begin{enumerate}
	\UCpaso[\UCactor] Gestiona los proyectos a través de las botones: \raisebox{-1mm}{\includegraphics[height=11pt]{images/Iconos/entrar}}, \raisebox{-1mm}{\includegraphics[height=11pt]{images/Iconos/pp}} y \raisebox{-1mm}{\includegraphics[height=11pt]{images/Iconos/word}}. \label{CU4-TB1}
	\item[- -] - - {\em {Fin del caso de uso}}.%
\end{enumerate}

\subsubsection{Puntos de extensión}
\UCExtenssionPoint{El actor requiere gestionar los módulos de un proyecto.}{Presionando el icono \raisebox{-1mm}{\includegraphics[height=11pt]{images/Iconos/entrar}} del paso \ref{CU4-P3} de la trayectoria principal o del paso \ref{CU4-TB1} de la trayectoria alternativa B.}{\UCref{CU5}{Gestionar Módulos}}
\UCExtenssionPoint{El actor requiere descargar el documento de análisis de un proyecto.}{Presionando el icono \raisebox{-1mm}{\includegraphics[height=11pt]{images/Iconos/pp}} o \raisebox{-1mm}{\includegraphics[height=11pt]{images/Iconos/word}} del paso \ref{CU4-P3} de la trayectoria principal o del paso \ref{CU4-TB1} de la trayectoria alternativa B.}{\UCref{CU13}{Descargar Documento}}
\UCExtenssionPoint{El actor requiere elegir a los colaboradores de un proyecto.}{Presionando el icono \raisebox{-1mm}{\includegraphics[height=11pt]{images/Iconos/Colaboradores}} del paso \ref{CU4-P3} de la trayectoria principal.}{\UCref{CU4.1}{Elegir Colaboradores}}
